% Options for packages loaded elsewhere
% Options for packages loaded elsewhere
\PassOptionsToPackage{unicode}{hyperref}
\PassOptionsToPackage{hyphens}{url}
\PassOptionsToPackage{dvipsnames,svgnames,x11names}{xcolor}
%
\documentclass[
  german,
  10pt,
  twoside]{article}
\usepackage{xcolor}
\usepackage[top=2.5cm,bottom=2.5cm,left=2cm,right=2cm]{geometry}
\usepackage{amsmath,amssymb}
\setcounter{secnumdepth}{5}
\usepackage{iftex}
\ifPDFTeX
  \usepackage[T1]{fontenc}
  \usepackage[utf8]{inputenc}
  \usepackage{textcomp} % provide euro and other symbols
\else % if luatex or xetex
  \usepackage{unicode-math} % this also loads fontspec
  \defaultfontfeatures{Scale=MatchLowercase}
  \defaultfontfeatures[\rmfamily]{Ligatures=TeX,Scale=1}
\fi
\usepackage{lmodern}
\ifPDFTeX\else
  % xetex/luatex font selection
\fi
% Use upquote if available, for straight quotes in verbatim environments
\IfFileExists{upquote.sty}{\usepackage{upquote}}{}
\IfFileExists{microtype.sty}{% use microtype if available
  \usepackage[]{microtype}
  \UseMicrotypeSet[protrusion]{basicmath} % disable protrusion for tt fonts
}{}
\usepackage{setspace}
\makeatletter
\@ifundefined{KOMAClassName}{% if non-KOMA class
  \IfFileExists{parskip.sty}{%
    \usepackage{parskip}
  }{% else
    \setlength{\parindent}{0pt}
    \setlength{\parskip}{6pt plus 2pt minus 1pt}}
}{% if KOMA class
  \KOMAoptions{parskip=half}}
\makeatother
% Make \paragraph and \subparagraph free-standing
\makeatletter
\ifx\paragraph\undefined\else
  \let\oldparagraph\paragraph
  \renewcommand{\paragraph}{
    \@ifstar
      \xxxParagraphStar
      \xxxParagraphNoStar
  }
  \newcommand{\xxxParagraphStar}[1]{\oldparagraph*{#1}\mbox{}}
  \newcommand{\xxxParagraphNoStar}[1]{\oldparagraph{#1}\mbox{}}
\fi
\ifx\subparagraph\undefined\else
  \let\oldsubparagraph\subparagraph
  \renewcommand{\subparagraph}{
    \@ifstar
      \xxxSubParagraphStar
      \xxxSubParagraphNoStar
  }
  \newcommand{\xxxSubParagraphStar}[1]{\oldsubparagraph*{#1}\mbox{}}
  \newcommand{\xxxSubParagraphNoStar}[1]{\oldsubparagraph{#1}\mbox{}}
\fi
\makeatother


\usepackage{graphicx}
\makeatletter
\newsavebox\pandoc@box
\newcommand*\pandocbounded[1]{% scales image to fit in text height/width
  \sbox\pandoc@box{#1}%
  \Gscale@div\@tempa{\textheight}{\dimexpr\ht\pandoc@box+\dp\pandoc@box\relax}%
  \Gscale@div\@tempb{\linewidth}{\wd\pandoc@box}%
  \ifdim\@tempb\p@<\@tempa\p@\let\@tempa\@tempb\fi% select the smaller of both
  \ifdim\@tempa\p@<\p@\scalebox{\@tempa}{\usebox\pandoc@box}%
  \else\usebox{\pandoc@box}%
  \fi%
}
% Set default figure placement to htbp
\def\fps@figure{htbp}
\makeatother



\ifLuaTeX
\usepackage[bidi=basic]{babel}
\else
\usepackage[bidi=default]{babel}
\fi
% get rid of language-specific shorthands (see #6817):
\let\LanguageShortHands\languageshorthands
\def\languageshorthands#1{}
\ifLuaTeX
  \usepackage[german]{selnolig} % disable illegal ligatures
\fi


\setlength{\emergencystretch}{3em} % prevent overfull lines

\providecommand{\tightlist}{%
  \setlength{\itemsep}{0pt}\setlength{\parskip}{0pt}}



 
\usepackage[style=authoryear,backend=biber,style=authoryear-comp,uniquename=false,uniquelist=false]{biblatex}
\addbibresource{../references.bib}


% Font configuration (Statis Sans alternative - use sans-serif throughout)
\renewcommand{\familydefault}{\sfdefault}
\usepackage{helvet}  % Use Helvetica-like font

\usepackage{csquotes}
\usepackage{booktabs}
\usepackage{array}
\usepackage{multirow}
\usepackage{wrapfig}
\usepackage{float}
\usepackage{colortbl}
\usepackage{pdflscape}
\usepackage{threeparttable}
\usepackage{threeparttablex}
\usepackage[normalem]{ulem}
\usepackage{makecell}
\usepackage{xcolor}
\usepackage{multicol}
\usepackage{fancyhdr}
\usepackage{titlesec}
\usepackage{etoolbox}
\usepackage{tabularx}
\usepackage{afterpage}
\usepackage{tikz}
\usepackage{longtable}
\usepackage{totcount}
\usepackage{xstring}

% Time display with Berlin timezone (+1 hour offset)
\newcount\berlinhr
\newcount\berlinmin
\newcommand{\berlintime}{%
  \berlinhr=\time\relax
  \divide\berlinhr by 60\relax
  \berlinmin=\time\relax
  \multiply\berlinhr by -60\relax
  \advance\berlinmin by \berlinhr\relax
  \berlinhr=\time\relax
  \divide\berlinhr by 60\relax
  \advance\berlinhr by 1\relax
  \ifnum\berlinhr>23\relax
    \advance\berlinhr by -24\relax
  \fi
  \ifnum\berlinhr<10 0\fi\the\berlinhr:\ifnum\berlinmin<10 0\fi\the\berlinmin%
}

% Register counters for total counting
\regtotcounter{table}
\regtotcounter{figure}

% Suppress default maketitle
\renewcommand{\maketitle}{}

% ===== DRAFT MODE CONTROL =====
% Set \draftmodetrue to enable draft mode (metadata page + ENTWURF header)
% Set \draftmodefalse to disable draft mode (for final version)
\newif\ifdraftmode
\draftmodetrue  % <-- Change to \draftmodefalse to disable draft mode

% Define default commands for counts (will be overridden by filter)
\providecommand{\totalchars}{0}
\providecommand{\totalwords}{0}
\providecommand{\abstractchars}{0}
\providecommand{\abstractwords}{0}
\providecommand{\zusammenfassungchars}{0}
\providecommand{\zusammenfassungwords}{0}
\providecommand{\footnotecountnum}{0}
\providecommand{\citationcountnum}{0}
\providecommand{\schlusselwortercount}{0}
\providecommand{\keywordscount}{0}

% Define metadata commands from Pandoc variables
\newcommand{\doctitle}{$title$}
\newcommand{\docauthors}{$author-short$}
\newcommand{\docpublication}{$publication$}
\newcommand{\docpublicationissue}{$publication-issue$}
\newcommand{\docshorttitle}{$short-title$}

\ifdraftmode
% Define metadata page command
\newcommand{\metadatapage}{%
  \clearpage
  \thispagestyle{empty}
  \onecolumn
  \vspace*{1.5cm}
  
  % Header section
  {\raggedright
  {\LARGE\bfseries\color{wistaBlue} Dokument-Metadaten}\\[0.5cm]
  
  \begin{tabular}{@{}ll}
    \textbf{Titel:} & \doctitle \\[0.3em]
    \textbf{Autoren:} & \docauthors \\[0.3em]
    \textbf{Veröffentlichung:} & \docpublication\ \docpublicationissue \\[0.3em]
    \textbf{Erstellungsdatum:} & \today\ um \berlintime\ Uhr \\[0.3em]
  \end{tabular}
  
  \vspace{0.5cm}
  {\color{wistaBlue}\rule{\textwidth}{1pt}}
  \par}
  
  \vspace{1cm}
  
  % Dokumentstatistik
  {\large\bfseries\color{wistaDarkGray} Dokumentstatistik}\\[0.5cm]
  
  \small
  \begin{tabular}{@{}lrrrrrr@{}}
    \toprule
    \textbf{Abschnitt} & \textbf{Wörter} & \textbf{Zeichen} & \textbf{Fußn.} & \textbf{Quellen} & \textbf{Vorgabe} \\
    \midrule
    Insgesamt & \textcolor{wistaBlue}{\textbf{\totalwords}} & \textcolor{wistaBlue}{\textbf{\totalchars}} & \textcolor{wistaBlue}{\textbf{\footnotecountnum}} & \textcolor{wistaBlue}{\textbf{\citationcountnum}} & 30\,000 Z. \\
    Zusammenfassung & \textcolor{wistaBlue}{\textbf{\zusammenfassungwords}} & \textcolor{wistaBlue}{\textbf{\zusammenfassungchars}} & --- & --- & 800 Z. \\
    Abstract & \textcolor{wistaBlue}{\textbf{\abstractwords}} & \textcolor{wistaBlue}{\textbf{\abstractchars}} & --- & --- & 700 Z. \\
    \midrule
    \sectionrows
    \bottomrule
  \end{tabular}
  \normalsize
  
  \vspace{0.5cm}
  
  % Keywords list
  {\large\bfseries\color{wistaDarkGray} Schlüsselwörter}\\[0.3cm]
  \small
  \schlusselworterlist
  \normalsize
  
  \vspace{0.5cm}
  
  {\large\bfseries\color{wistaDarkGray} Keywords}\\[0.3cm]
  \small
  \keywordslist
  \normalsize
  
  \vspace{1cm}
  
  % Tables list
  {\large\bfseries\color{wistaDarkGray} Tabellen}\\[0.3cm]
  
  \small
  \begin{tabular}{@{}lp{12cm}@{}}
    \toprule
    \textbf{Nr.} & \textbf{Beschreibung} \\
    \midrule
    \tablerows
    \bottomrule
  \end{tabular}
  \normalsize
  
  \vspace{1cm}
  
  % Figures list
  {\large\bfseries\color{wistaDarkGray} Abbildungen}\\[0.3cm]
  
  \small
  \begin{tabular}{@{}lp{12cm}@{}}
    \toprule
    \textbf{Nr.} & \textbf{Beschreibung} \\
    \midrule
    \figurerows
    \bottomrule
  \end{tabular}
  \normalsize
  
  \clearpage
}
\fi

% Define WISTA colors (from Kristiansen Word template)
\definecolor{wistaBlue}{RGB}{91,155,213}
\definecolor{wistaLightBlue}{RGB}{120,170,215}
\definecolor{wistaDarkGray}{RGB}{96,94,92}
\definecolor{wistaLightGray}{RGB}{128,128,128}

% Section formatting (##) - number 20% bigger, lines closer, text closer to bottom line
% e.g., "2" with lines around "Einleitung"
\titleformat{\section}[display]
  {\normalfont\large\bfseries}
  {\color{wistaBlue}\Large\bfseries\thesection\\\noindent\rule{\columnwidth}{1.5pt}}
  {0pt}
  {\vspace{0pt}\color{wistaDarkGray}\large\bfseries}
  [\vspace{0pt}\noindent{\color{wistaBlue}\rule{\columnwidth}{1.5pt}}\vspace{0.05cm}]
\titlespacing*{\section}{0pt}{1.5ex plus 0.3ex minus .1ex}{0.2ex plus .1ex}

% Subsection formatting (###) - grey text with blue line under
% e.g., "2.1 Handlungsbedarf"
\renewcommand{\thesubsection}{\thesection.\arabic{subsection}}
\titleformat{\subsection}
  {\color{wistaDarkGray}\normalsize\bfseries}
  {\color{wistaDarkGray}\normalsize\bfseries\thesubsection}
  {0.5em}
  {\color{wistaDarkGray}\normalsize\bfseries}
  [\vspace{0pt}\noindent{\color{wistaBlue}\rule{\columnwidth}{0.8pt}}\vspace{0.05cm}]
\titlespacing*{\subsection}{0pt}{1ex plus 0.2ex minus .1ex}{0.2ex plus .1ex}
  
% Subsubsection formatting (####) - grey text, no underline
% e.g., "2.1.1 Steigende Datenmenge"
\renewcommand{\thesubsubsection}{\thesubsection.\arabic{subsubsection}}
\titleformat{\subsubsection}
  {\color{wistaDarkGray}\normalsize\bfseries}
  {\color{wistaDarkGray}\normalsize\bfseries\thesubsubsection}
  {0.5em}
  {\color{wistaDarkGray}\normalsize\bfseries}

% Paragraph formatting (4th level) - numbered as subsection.subsubsection.paragraph
\renewcommand{\theparagraph}{\thesubsubsection.\arabic{paragraph}}
\titleformat{\paragraph}[runin]
  {\color{wistaDarkGray}\normalsize\bfseries}
  {\color{wistaBlue}\normalsize\bfseries\theparagraph}
  {0.5em}
  {\color{wistaDarkGray}\normalsize\bfseries}
  [.]
\titlespacing*{\paragraph}{0pt}{1ex plus 0.5ex minus 0.2ex}{0.5em}

% Header and footer setup with blue line under header
\pagestyle{fancy}
\fancyhf{}
\fancyhead[L]{\small\textit{\docauthors}}
\ifdraftmode
  \fancyhead[C]{\small\color{red}\textbf{ENTWURF} -- \today}
\fi
\fancyhead[R]{\small\textit{\docshorttitle}}
\fancyfoot[L]{\small\thepage}
\fancyfoot[R]{\small Statistisches Bundesamt | \docpublication\ | {\color{wistaBlue}\docpublicationissue}}
\renewcommand{\headrulewidth}{0pt}
\renewcommand{\footrulewidth}{0pt}
% Add blue line under header
\renewcommand{\headrule}{\color{wistaBlue}\hrule width\headwidth height 1pt \vskip 2pt}
\renewcommand{\headrulewidth}{1pt}

% Make all figures span both columns (full width)
\let\origfigure\figure
\let\endorigfigure\endfigure
\renewenvironment{figure}[1][htbp]{%
  \origfigure*[#1]%
}{%
  \endorigfigure*%
}

% Allow manual control of table* environment for full-width tables
% Tables by default stay in single column, but can use table* manually



% Make bibliography full-width (one column) - for biblatex
\defbibheading{bibliography}[\bibname]{%
  \onecolumn
  \section*{#1}%
  \markboth{#1}{#1}%
}



% Footnotes stay within column
\usepackage[stable]{footmisc}

% Custom footnote formatting: "|^1" in dark WISTA gray
\renewcommand{\thefootnote}{{\color{wistaDarkGray}|$^{\arabic{footnote}}$}}

% Pifont for dingbat symbols
\usepackage{pifont}

% Better table handling for two-column mode
\usepackage{array,tabularx,calc}

% Custom WISTA table format
\usepackage{caption}
\usepackage{xcolor}
\usepackage{colortbl}
\usepackage{array}

% Define custom caption format for WISTA tables
\DeclareCaptionFormat{wista}{%
  {\color{wistaLightBlue}\bfseries #1#2}\par\vspace{0.2em}%
  {\color{wistaDarkGray}\normalfont #3}\par\vspace{0.3em}%
}

% Set up table captions with custom WISTA format
\captionsetup[table]{
  position=top,
  format=wista,
  justification=justified,
  singlelinecheck=false,
  labelsep=none,
  skip=0.3em
}

% Better table spacing and white inner lines
\setlength{\tabcolsep}{6pt}
\renewcommand{\arraystretch}{1.15}

% Set white lines for tables by default, blue for first column
\AtBeginDocument{
  \arrayrulecolor{white}
  \setlength{\arrayrulewidth}{0.3pt}
}

% Define commands for colored table borders
\newcommand{\bluevrule}{\color{wistaBlue}\vrule}
\newcommand{\whitevrule}{\color{white}\vrule}

% Define WISTA table colors
\definecolor{wistaTableBg}{RGB}{240,248,255}     % Very light blue background
\definecolor{wistaHeaderBg}{RGB}{220,230,240}    % Light blue header background
\definecolor{wistaFirstCol}{RGB}{65,125,180}     % Blue for first column lines

% Custom WISTA table command
\newcommand{\wistaTableStart}{%
  \rowcolors{2}{wistaTableBg}{wistaTableBg}%
  \arrayrulecolor{wistaFirstCol}%
}

% Custom commands for table styling
\newcommand{\wistaHeaderRow}[1]{%
  \rowcolor{wistaHeaderBg}%
  \arrayrulecolor{wistaDarkGray}%
  #1 \\
  \arrayrulecolor{white}%
}

\newcommand{\wistaFirstColRule}{\arrayrulecolor{wistaFirstCol}}
\newcommand{\wistaWhiteRule}{\arrayrulecolor{white}}

% Custom bullet points with > character - 35% width, 95% height, black extra bold, half line spacing
\usepackage{enumitem}
\setlist[itemize]{label={\scalebox{0.35}[0.95]{\fontseries{bx}\selectfont >}}, leftmargin=1.5em, itemsep=5pt, topsep=3pt, parsep=0pt}

% Redefine ref to include pifont 247 (arrow) for figure/table/section references
\let\oldref\ref
\renewcommand{\ref}[1]{\oldref{#1}\raisebox{-0.2ex}{\scriptsize\color{wistaBlue}\ding{247}}}
\makeatletter
\@ifpackageloaded{caption}{}{\usepackage{caption}}
\AtBeginDocument{%
\ifdefined\contentsname
  \renewcommand*\contentsname{Inhaltsverzeichnis}
\else
  \newcommand\contentsname{Inhaltsverzeichnis}
\fi
\ifdefined\listfigurename
  \renewcommand*\listfigurename{Abbildungsverzeichnis}
\else
  \newcommand\listfigurename{Abbildungsverzeichnis}
\fi
\ifdefined\listtablename
  \renewcommand*\listtablename{Tabellenverzeichnis}
\else
  \newcommand\listtablename{Tabellenverzeichnis}
\fi
\ifdefined\figurename
  \renewcommand*\figurename{Abbildung}
\else
  \newcommand\figurename{Abbildung}
\fi
\ifdefined\tablename
  \renewcommand*\tablename{Tabelle}
\else
  \newcommand\tablename{Tabelle}
\fi
}
\@ifpackageloaded{float}{}{\usepackage{float}}
\floatstyle{ruled}
\@ifundefined{c@chapter}{\newfloat{codelisting}{h}{lop}}{\newfloat{codelisting}{h}{lop}[chapter]}
\floatname{codelisting}{Listing}
\newcommand*\listoflistings{\listof{codelisting}{Listingverzeichnis}}
\makeatother
\makeatletter
\makeatother
\makeatletter
\@ifpackageloaded{caption}{}{\usepackage{caption}}
\@ifpackageloaded{subcaption}{}{\usepackage{subcaption}}
\makeatother
\usepackage{bookmark}
\IfFileExists{xurl.sty}{\usepackage{xurl}}{} % add URL line breaks if available
\urlstyle{same}
\hypersetup{
  pdftitle={EFFIZIENTE AUSWERTUNG DES BUSINESS-TAX-PANELS MITHILFE VON R},
  pdfauthor={Oliver Hauke; Annette Kristiansen},
  pdflang={de},
  colorlinks=true,
  linkcolor={blue},
  filecolor={Maroon},
  citecolor={Blue},
  urlcolor={Blue},
  pdfcreator={LaTeX via pandoc}}


\title{EFFIZIENTE AUSWERTUNG DES BUSINESS-TAX-PANELS MITHILFE VON R}
\author{Oliver Hauke \and Annette Kristiansen}
\date{2025-12-05}
\begin{document}
\maketitle

% Custom WISTA-style title page
\begin{titlepage}
\thispagestyle{empty}
\pagestyle{empty}
\onecolumn

% Horizontal blue line at top (header line)
\noindent{\color{wistaBlue}\rule{\textwidth}{1pt}}

\vspace{-0.5cm}

% Vertical blue line on the left (6.3cm from left edge, starts below header line, stops above footer)
\begin{tikzpicture}[remember picture, overlay]
  \draw[wistaBlue, line width=2pt] 
    ([xshift=6.3cm, yshift=-4cm]current page.north west) -- 
    ([xshift=6.3cm, yshift=3.5cm]current page.south west);
\end{tikzpicture}

\vspace*{0.5cm}

% Left side - short author bios (right-aligned to the vertical line)
\begin{minipage}[t]{3.9cm}
\selectlanguage{ngerman}%
\raggedleft
\hyphenpenalty=0
\tolerance=1000
\fontsize{7}{8.5}\selectfont
\vspace{0pt}%
{\fontsize{8}{9.5}\selectfont\color{wistaBlue}\textbf{Oliver Hauke}}\\[0.15cm]
ist Mathematiker und IT-Referent im Kompetenzzentrum „Analyse und Auswertung" des Statistischen Bundesamtes. Er verantwortet die Weiterentwicklung der technischen Infrastruktur des Forschungsdatenzentrums und berät das Netzwerk empirische Steuerforschung in der effizienten Nutzung der Systeme.\\[0.5cm]

{\fontsize{8}{9.5}\selectfont\color{wistaBlue}\textbf{Annette Kristiansen}}\\[0.15cm]
ist Ökonomin und Referentin im Referat „Unternehmenssteuern" des Statistischen Bundesamtes. Sie beschäftigt sich mit der Zusammenführung der Unternehmenssteuerstatistiken und betreut das Netzwerk empirische Steuerforschung. Für ihre externe Promotion an der Universität Bremen untersucht sie die technologische Vielfalt multinationaler Unternehmen.

\end{minipage}%
\hspace{0.4cm}% Gap to vertical line (left side)
\hspace{0.4cm}% Gap from vertical line (right side)
% Right side - title, authors, keywords, abstracts
\begin{minipage}[t]{11.8cm}
\raggedright
\vspace{1.2cm}%

% Title (uppercase via command, not bold)
{\LARGE\color{wistaDarkGray}\begin{spacing}{1.3}\MakeUppercase{Effiziente Auswertung des Business-Tax-Panels mithilfe von R}\end{spacing}}

\vspace{-0.2cm}

% Blue horizontal line (underline for title)
{\color{wistaBlue}\rule{\linewidth}{0.5pt}}

\vspace{0.5cm}

% Authors inline, dark gray
{\large\color{wistaDarkGray} Oliver Hauke, Annette Kristiansen}

\vspace{0.05cm}

% Blue line (underline for authors)
{\color{wistaBlue}\rule{\linewidth}{0.5pt}}

\vspace{0.5cm}

% Keywords (pifont 247 arrow then bold label and blue text)
\hangindent=1.2em\hangafter=1
{\color{wistaBlue}\ding{247}\ {\bfseries Schlüsselwörter:} Effiziente Analysen, R-Programmierung, Forschungsdaten, Unternehmenssteuerstatistik, Business-Tax-Panel, Netzwerk empirische Steuerforschung}

\vspace{0.3cm}

% Zusammenfassung (uppercase via command)
{\color{wistaDarkGray}\bfseries\MakeUppercase{Zusammenfassung}}\\[0.15cm]
\small
Das Forschungsdatenzentrum des Statistischen Bundesamtes baut seine technische Infrastruktur gezielt aus, um der Wissenschaft erweiterte Analysemöglichkeiten amtlicher Daten bereitzustellen. Seit November 2025 steht Forschenden exklusiver Zugang zu diesen erweiterten Systemressourcen in R und Stata zur Verfügung. R-basierte Tools -- insbesondere Apache Arrow, Parquet und DuckDB -- ermöglichen es, gängige Analysen großer Forschungsdatensätze in Echtzeit auszuführen. Am Beispiel des Business-Tax-Panel (2013 bis 2019) werden die erzielbaren Effizienzgewinne bei der Auswertung umfangreicher steuerstatistischer Daten im Rahmen des Netzwerks empirische Steuerforschung demonstriert.

\vspace{0.5cm}

% English keywords and abstract (italicized)
\hangindent=1.2em\hangafter=1
{\itshape\color{wistaBlue}\ding{247}\ {\bfseries Keywords:} Efficient analysis -- R programming -- research data -- corporate tax statistics -- Business Tax Panel -- empirical tax research network}

\vspace{0.3cm}

{\itshape\color{wistaDarkGray}\bfseries\MakeUppercase{Abstract}}\\[0.15cm]
\small\itshape
The Research Data Centre at the Federal Statistical Office is strategically expanding its technical infrastructure to provide researchers with enhanced analytical capabilities for official data. Since November 2025, researchers have had exclusive access to these extended system resources in R and Stata. R-based tools -- particularly Apache Arrow, Parquet, and DuckDB -- enable real-time analysis of large research datasets. Using the Business Tax Panel (2013--2019), this paper demonstrates the achievable efficiency gains in analysing large-scale tax statistics within the empirical tax research network.

\end{minipage}

\vfill

\vspace{0.5cm}
% Footer matching document style
\noindent\makebox[\textwidth]{%
  \small 1%
  \hfill%
  \small Statistisches Bundesamt | WISTA | {\color{wistaBlue}01/2026}%
}

\end{titlepage}

% Stay in one-column mode for TOC (will be switched later)
\pagestyle{fancy}

\renewcommand*\contentsname{Inhaltsverzeichnis}
{
\hypersetup{linkcolor=}
\setcounter{tocdepth}{3}
\tableofcontents
}

\setstretch{1}
\renewcommand{\totalchars}{6 256}
\renewcommand{\totalwords}{714}
\renewcommand{\abstractchars}{0}
\renewcommand{\abstractwords}{0}
\renewcommand{\zusammenfassungchars}{0}
\renewcommand{\zusammenfassungwords}{0}
\renewcommand{\footnotecountnum}{1}
\renewcommand{\citationcountnum}{0}
\renewcommand{\schlusselwortercount}{0}
\renewcommand{\keywordscount}{0}
\renewcommand{\doctitle}{EFFIZIENTE AUSWERTUNG DES BUSINESS-TAX-PANELS MITHILFE VON R}
\renewcommand{\docauthors}{Oliver Hauke, Annette Kristiansen}
\renewcommand{\docpublication}{WISTA}
\renewcommand{\docpublicationissue}{01/2026}
\renewcommand{\docshorttitle}{Effiziente Analyse großer Forschungsdaten}
\newcommand{\sectionrows}{            Einleitung & \textcolor{wistaBlue}{\textbf{1 231}} & \textcolor{wistaBlue}{\textbf{2 026}} & \textcolor{wistaBlue}{\textbf{0}} & \textcolor{wistaBlue}{\textbf{0}} & --- \\
            Anwendungsfall Business-Tax-Panel & \textcolor{wistaBlue}{\textbf{125}} & \textcolor{wistaBlue}{\textbf{205}} & \textcolor{wistaBlue}{\textbf{0}} & \textcolor{wistaBlue}{\textbf{0}} & --- \\
            Techniche Infrastruktur & \textcolor{wistaBlue}{\textbf{145}} & \textcolor{wistaBlue}{\textbf{238}} & \textcolor{wistaBlue}{\textbf{0}} & \textcolor{wistaBlue}{\textbf{0}} & --- \\
            Datenverarbeitungsmethoden & \textcolor{wistaBlue}{\textbf{810}} & \textcolor{wistaBlue}{\textbf{1 333}} & \textcolor{wistaBlue}{\textbf{0}} & \textcolor{wistaBlue}{\textbf{0}} & --- \\
            Performanzvergleich & \textcolor{wistaBlue}{\textbf{999}} & \textcolor{wistaBlue}{\textbf{1 644}} & \textcolor{wistaBlue}{\textbf{0}} & \textcolor{wistaBlue}{\textbf{0}} & --- \\
            Ergebnis für das vollständige BTP & \textcolor{wistaBlue}{\textbf{84}} & \textcolor{wistaBlue}{\textbf{138}} & \textcolor{wistaBlue}{\textbf{0}} & \textcolor{wistaBlue}{\textbf{0}} & --- \\
            Fazit & \textcolor{wistaBlue}{\textbf{407}} & \textcolor{wistaBlue}{\textbf{669}} & \textcolor{wistaBlue}{\textbf{0}} & \textcolor{wistaBlue}{\textbf{0}} & --- \\
}
\newcommand{\tablerows}{            \hyperref[tbl-pkg-versions]{Tabelle~1} & Versionen verfügbarer R-Packages \\
            \hyperref[tbl-baseline]{Tabelle~2} & Baselines in SAS und Stata (anhand 1\% des simuliertes BTP) \\
            \hyperref[tbl-benchmark-small]{Tabelle~3} & Performanz Messungen für die Fallzahlberechnung anhand 1\% des BTP \\
            \hyperref[tbl-benchmark-regr-small]{Tabelle~4} & Performanz Messungen für die Regressionsberechnung anhand 1\% des BTP \\
            \hyperref[tbl-benchmark]{Tabelle~5} & Performanz Messungen für die Fallzahlberechnung anhand 10\% des BTP \\
            \hyperref[tbl-benchmark-regr]{Tabelle~6} & Performanz Messungen für die Regressionsberechnung anhand 10\% des BTP \\
            \hyperref[tbl-bm-cnt-small]{Tabelle~7} & Benchmarkergebnisse für Fallzahlen (kleine Stichproben des BTP) \\
            \hyperref[tbl-bm-reg-small]{Tabelle~8} & Benchmarkergebnisse für Regression (kleine Stichproben des BTP) \\
}
\newcommand{\figurerows}{            \hyperref[fig-optimization-strategies]{Abbildung~1} & Strategien zur Performanz-Optimierung im FDZ-Bund \\
}
\newcommand{\schlusselworterlist}{}
\newcommand{\keywordslist}{}
\ifdraftmode
\metadatapage
\fi

\section{Anwendungsfall
Business-Tax-Panel}\label{anwendungsfall-business-tax-panel}

\subsection{Datengrundlage}\label{datengrundlage}

\emph{Das Business-Tax-Panel kombiniert die Angaben aus verschiedenen
Ertrags- sowie Umsatzsteuerstatistiken im Quer- und Längsschnitt.
Verfügbar sind Informationen aus den folgenden Statistiken: Gewerbe-,
Körperschaftsteuer, Umsatzsteuer-Voranmeldung und -Veranlagung,
Statistik über die Personengesellschaften und Gemeinschaften,
Einnahmenüberschussrechnung und einige Merkmale aus dem Statistischen
Unternehmensregister. Insgesamt deckt der Forschungsdatensatz eine
Vielzahl an Analysemöglichkeiten, wie zum Beispiel Analyse von
Anpassungsreaktionen oder Verteilungsanalysen, die für die
evidenzbasierte steuerpolitische Entscheidungen nötig sind, ab.}

\emph{Für die vielfältigen Analysewünsche der Forschenden, der damit
einhergenden Dokumentation des Datensatzes in den Metadatenreports
(!!hier Fußnote für MDR - siehe Deskriptive Ecktaten 1.) des FDZ, und
den Auswertungsbedarf des Fachbereiches ist der Datensatz möglichst
wenig beschränkt worden und dadurch ressourcenintensive. Anstelle einer
Stichprobe wurden die zugrundeliegenden Vollerhebungen verwendet. Die
erste Version des BTP (2013 bis 2019) besteht aus ca. 67 Millionen
Beobachtungen. Mit der Aktualisierung um das Berichtsjahr 2020 wird die
Zahl der Beobachtungen sich auf ca. 78 Millionen erhöhen. Auch der
Merkmalsumfang von zurzeit über 2.700 Merkmalen wird durch Erweiterung
von Berichtsjahren steigen.}

\emph{Bereits aktuell und insbesondere durch die regelmäßige Erweiterung
um Bereichtsjahre steht der Fachbereich sowie das FDZ vor der
Herausforderung den hohe Bedarf an Speicherkapazität, Arbeitsspeicher
und Prozessorleistung zu decken. Daher wird das BTP als Anwendungsfall
herangezogen}

\subsection{Auswertungsbedarfe}\label{auswertungsbedarfe}

\emph{Um den Auswetungsbedarf des Fachbereichs und des FDZ abzubilden
wird die Tabelle zur Häufigkeitsverteilung der Beobachtungen je Jahr und
Steuerstatistik (FDZ, 2024 --\textgreater{} IN QULLEN aufnehmen!)
verwendet. Im Fachbereich ergibt sich zusätzlich der Bedarf die Daten je
nach Statistik und Tabelle zu bearbeiten. Auch die Forschenden
bearbeiten die Daten um sie für die jeweilige Analyse aufzubereiten.
Daher wird die Datenbearbeitung und anschließende Analyse als zweiter
Anwendungsbedarf getestet.}

\emph{Die Körperschaftsteuerstatistik zeigt über die Jahre eine Zunahme
der Steuerpflichtigen mit Verlustvortrag im jeweiligen Berichtsjahr
(GENESIS Code 73211-0002). Für den zweiten Anwendungsfall wird das
Vorliegen eines Verlustvortrages im jeweiligen Berichtsjahr als versucht
anhand verschiedener Einflussfaktoren zu erklären. Anzunehmen ist, dass
Personalkosten den Verlustvortrag erhöhen können, daher wird die Zahl
der Sozialversicherungspflichtig Beschäftigen sowie die Summe der Löhne
und Gehälter als erklärende Variable verwendet. Der jeweilige
Gesamtbetrag der Einkünfte und die abgezogenen Spenden können ebenfalls
den Verlustvortrag beeinflussen. Um den Internationalisierungsgrad des
Unternehmens zu erfassen wird ein Dummy für aufgenommen der Anzeigt, ob
das Unternehmen auslandskontrolliert ist.}

\subsubsection{Deskriptive Eckdaten}\label{deskriptive-eckdaten}

\begin{enumerate}
\def\labelenumi{\arabic{enumi}.}
\tightlist
\item
  deskriptive Eckdaten, für BTP als Panel im Quer- und Längsschnitt:
  Fallzahl pro Statistik und Jahr (s.MDR Produkt S. 9f
  https://www.forschungsdatenzentrum.de/sites/default/files/btp\_2013-2019\_on-site\_mdr-prod.pdf)
\end{enumerate}

Kennzahlen um inhaltlose Testdaten zu generieren werden folgende
Eckdaten benötigt und im nächsten BTP berücksichtigt:

\begin{itemize}
\tightlist
\item
  Befüllung
\item
  Statistische Eckdaten univariate (Quantile),
\item
  ggf. multivariate (gemeinsame Verteilung)
\end{itemize}

Exemplarisch, neues BTP beinhaltet dies. Regelmäßige Updates
beschleunigen.

\subsubsection{Regressionen}\label{regressionen}

\begin{enumerate}
\def\labelenumi{\arabic{enumi}.}
\setcounter{enumi}{1}
\tightlist
\item
  Verlustvorträge\footnote{Vortrag Annette 2024 JK} (s. Vortrag NeSt)
  als Indikator für \ldots{} und Einflussgrößen auf Verlustvorträge, für
  Vermutung von Einfluss der Unternehmensgröße (großes Gesamtvolumen,
  Löhne, Spenden, tätige Personen; Auslandskontrolle oder WZ) auf
  positive Verlustvorträge.
\end{enumerate}

\subsection{Technische
Spezifikationen}\label{technische-spezifikationen}

\emph{AK: Habe ich im Bereich Datengrundlage integriert. }

\emph{Für die vielfältigen Analysewünsche der Forschenden, den
Auswertungsbedarf des Fachbereiches bspw. für Sonderauswertungen und der
Erstellung des Metadatenreports des Forschungsdatenzentrums ist der
Forschungsdatensatz konzipiert. Für die Abdeckung aller Anforderungen
ist der Datensatz möglichst wenig beschränkt worden und dadurch
ressourcenintensive. Der damit einhergehende hohe Bedarf an
Speicherkapazität, Arbeitsspeicher und Prozessorleistung stellt sowohl
im Fachbereich als auch im Forschungsdatenzentrum eine Herausforderung
dar.}

\emph{Für die vielfältigen Analysewünsche der Forschenden, den
Auswertungsbedarf des Fachbereiches bspw. für Sonderauswertungen und der
Erstellung des Metadatenreports des Forschungsdatenzentrums ist der
Forschungsdatensatz konzipiert. Für die Abdeckung aller Anforderungen
ist der Datensatz möglichst wenig beschränkt worden und dadurch
ressourcenintensive. Der damit einhergehende hohe Bedarf an
Speicherkapazität, Arbeitsspeicher und Prozessorleistung stellt sowohl
im Fachbereich als auch im Forschungsdatenzentrum eine Herausforderung
dar.}

\begin{itemize}
\item
  Technische Herausforderung:

  \begin{itemize}
  \tightlist
  \item
    effiziente und verständliche Speicherung
  \item
    Herausfiltern großer Datenmengen irrelevanter Merkmale
  \end{itemize}
\item
  Umfang über verfügbaren Arbeitsspeicher
\item
  Besonders breit mit 3000 Merkmalen, 69 Mio. Beobachtung, sodass
  zeilenbasierte Formate problematisch sind.
\item
  Technische Herausforderung:

  \begin{itemize}
  \tightlist
  \item
    effiziente und verständliche Speicherung
  \item
    Herausfiltern großer Datenmengen irrelevanter Merkmale
  \end{itemize}
\end{itemize}


\printbibliography



\end{document}
